%! Matrices as Functions — Unit 4, Lesson 4.13 — Guided Notes
\documentclass[10pt]{article}
\usepackage{precalculus-article}
\usepackage{precalculus-boxes}
\newcommand{\TallMath}[1]{$\displaystyle #1\rule[-1.4em]{0pt}{3.2em}$}

\begin{document}

\pageheader{Unit 4, Lesson 4.13}{Guided Notes}
\namedateperiod

\vspace{0.08in}

\begin{objectivebox}
By the end of this lesson, I will be able to:
\begin{itemize}[leftmargin=1.5em, itemsep=3pt]
  \item \writeline
  \item \writeline
  \item \writeline
  \item \writeline
\end{itemize}
\end{objectivebox}

\vspace{0.06in}

\begin{vocabbox}
\termblanklong{linear transformation}
\termblanklong{transformation matrix}
\termblanklong{rotation matrix}
\termblanklong{dilation (area scale factor)}
\termblanklong{composition of transformations}
\termblanklong{inverse transformation}
\end{vocabbox}

\vspace{0.06in}

\begin{hookbox}
To rotate a logo $90°$, a graphic-design program applies the same matrix to every point in
the image. What matrix achieves a $90°$ counterclockwise rotation? If you rotate $30°$ then $60°$,
is that the same as rotating $90°$ all at once?

\writeline
\end{hookbox}

\vspace{0.06in}

\begin{notesbox}{Section 1 --- Reading the Matrix from Unit-Vector Images}
\small
Every linear transformation $L$ is determined by the images of the standard unit vectors
$\hat{\imath}=\langle1,0\rangle$ and $\hat{\jmath}=\langle0,1\rangle$.

\vspace{0.06in}
\textbf{Rule.} If $L(\hat{\imath})=\langle a_{11},a_{21}\rangle$ and
$L(\hat{\jmath})=\langle a_{12},a_{22}\rangle$, then the associated matrix is:
\[
  A = \begin{bmatrix} \blank{1.2cm} & \blank{1.2cm} \\ \blank{1.2cm} & \blank{1.2cm} \end{bmatrix}
\]

\vspace{0.04in}
\textbf{Key idea:} the \emph{columns} of $A$ are the images of \blank{2cm} and \blank{2cm}.

\vspace{0.08in}
\textbf{Example.} $L$ maps $\hat{\imath} \to \langle3,1\rangle$ and $\hat{\jmath} \to \langle-1,2\rangle$.

\hspace{1em} $A = \begin{bmatrix} \blank{1cm} & \blank{1cm} \\ \blank{1cm} & \blank{1cm} \end{bmatrix}$
\qquad Apply to $\langle2,-1\rangle$: \; $A\vec v = $ \blank{4cm}
\end{notesbox}

\vspace{0.06in}

\begin{notesbox}{Section 2 --- The Rotation Matrix}
\small
The matrix that rotates every vector counterclockwise by angle $\theta$:
\[
  R(\theta) = \begin{bmatrix} \blank{2cm} & \blank{2cm} \\ \blank{2cm} & \blank{2cm} \end{bmatrix}
\]

\textbf{Why?} $R(\theta)$ maps $\hat{\imath} \to $ \blank{3cm} and $\hat{\jmath} \to $ \blank{3cm}.

\vspace{0.08in}
\textbf{Example.} Write $R(90°)$ using exact values and apply it to $\langle3,1\rangle$.

\hspace{1em} $R(90°) = \begin{bmatrix} \blank{1cm} & \blank{1cm} \\ \blank{1cm} & \blank{1cm} \end{bmatrix}$
\qquad $R(90°)\begin{bmatrix}3\\1\end{bmatrix} = $ \blank{4cm}

\vspace{0.08in}
\textbf{Dilation (area scale).} $|\det A| = $ \blank{2cm} gives the factor by which $L$ scales areas.

\hspace{1em} For $A = \begin{bmatrix}2&0\\0&3\end{bmatrix}$: \; $\det A = $ \blank{1.2cm},
so every region's area is scaled by \blank{1.2cm}.
\end{notesbox}

\vspace{0.06in}

\begin{notesbox}{Section 3 --- Composing Transformations}
\small
Applying $L_1$ (matrix $A$) \textbf{first}, then $L_2$ (matrix $B$):

\[
  (L_2 \circ L_1)(\vec v) = B(A\vec v) = (BA)\vec v
  \qquad \Rightarrow \qquad \text{composition matrix} = \blank{1.5cm}
\]

\textbf{Order matters:} in general, $BA \ne AB$.

\vspace{0.08in}
\textbf{Example.} $A = R(30°)$, $B = R(60°)$. Claim: $BA = R(90°)$.

\hspace{1em}
$R(30°) = \begin{bmatrix}\frac{\sqrt{3}}{2} & -\frac{1}{2} \\[3pt] \frac{1}{2} & \frac{\sqrt{3}}{2}\end{bmatrix}$,\quad
$R(60°) = \begin{bmatrix}\frac{1}{2} & -\frac{\sqrt{3}}{2} \\[3pt] \frac{\sqrt{3}}{2} & \frac{1}{2}\end{bmatrix}$

\vspace{0.08in}
\hspace{1em} Compute $R(60°)\cdot R(30°)$: \writelines{3}
\end{notesbox}

\vspace{0.06in}

\begin{notesbox}{Section 4 --- Inverse Transformations}
\small
$L$ and $M$ are \textbf{inverse transformations} iff $(M \circ L)(\vec v) = \vec v$ for every $\vec v$,
i.e., the composition matrix equals \blank{1.2cm}.

\[
  \text{If } L(\vec v) = A\vec v, \text{ then } L^{-1}(\vec v) = \blank{3cm}.
\]

\textbf{Example.} $A = \begin{bmatrix}2&1\\1&1\end{bmatrix}$.

\hspace{1em} $\det A = $ \blank{1.5cm} \qquad
$A^{-1} = \dfrac{1}{\det A}\begin{bmatrix}\blank{0.8cm}&\blank{0.8cm}\\\blank{0.8cm}&\blank{0.8cm}\end{bmatrix}
= \begin{bmatrix}\blank{0.8cm}&\blank{0.8cm}\\\blank{0.8cm}&\blank{0.8cm}\end{bmatrix}$

\vspace{0.08in}
\hspace{1em} Verify $A\cdot A^{-1} = I$: \writelines{2}
\end{notesbox}

\vspace{0.06in}

\begin{practicebox}
\small
\textbf{Practice.} A transformation maps $\hat{\imath} \to \langle0,-1\rangle$ and $\hat{\jmath} \to \langle1,0\rangle$.

(a) Write matrix $A$: \blank{3cm} \qquad
(b) What rotation angle does this represent? \blank{2.5cm}

(c) Find $A^{-1}$: \blank{3cm} \qquad
(d) Describe $L^{-1}$ geometrically: \blank{5cm}
\end{practicebox}

\end{document}
